\section{NLP/ML 技术在语音 BCI 中的关键角色}

本节系统梳理语音 BCI 中使用的核心 NLP 和机器学习技术，帮助读者将 CS224N 的知识与 BCI 应用联系起来。

\subsection{序列到序列建模}

语音 BCI 的核心任务本质上是一个\textbf{序列到序列（seq2seq）}问题：

$$\text{神经特征序列} \xrightarrow{\text{Decoder}} \text{音素/词序列}$$

与 NMT（神经机器翻译）的对比：

\begin{table}[H]
\centering
\begin{tabular}{lcc}
\toprule
& \textbf{NMT} & \textbf{语音 BCI} \\
\midrule
输入 & 源语言 token 序列 & 神经特征向量序列 \\
输出 & 目标语言 token 序列 & 音素/词序列 \\
对齐方式 & 任意（可重排） & 单调（时间顺序） \\
数据量 & 百万级句对 & $\sim$10,000 句 \\
延迟要求 & 可容忍数百 ms & 必须 $<$ 20 ms \\
模型选择 & Transformer & GRU + CTC \\
\bottomrule
\end{tabular}
\caption{NMT 与语音 BCI 的 seq2seq 任务对比}
\end{table}

\begin{warningbox}{不要盲目照搬 NLP 模型}
虽然语音 BCI 在形式上是 seq2seq 问题，但其约束条件与典型 NLP 任务差异巨大：数据量小三个数量级、延迟要求严格百倍、对齐是单调的。直接使用 NLP 领域的``最强模型''（如 Transformer encoder-decoder）往往不是最优选择。\textbf{问题特性决定模型选择}，而非模型的``先进程度''。
\end{warningbox}

\subsection{语言模型的角色}

语言模型在语音 BCI 中扮演着\textbf{至关重要}的角色——它弥补了神经解码器的不足：

\begin{itemize}
    \item 神经解码器产生的音素序列可能包含错误（如将 /b/ 误认为 /p/）
    \item 语言模型知道``I can speak''比``I can speke''更可能，从而纠正解码错误
    \item N-gram 模型在 Beam Search 中实时引导搜索方向
    \item Transformer 语言模型在后处理中进行全局重排序
\end{itemize}

\begin{importantbox}{语言模型是``免费的先验知识''}
语言模型本质上编码了\textbf{语言的统计规律}——哪些词组合更常见、哪些句子更自然。这些知识不需要从患者的稀缺脑数据中学习，而是从海量文本语料中预训练获得的。在数据极度稀缺的 BCI 场景中，语言模型提供的先验知识尤为宝贵。
\end{importantbox}

\subsection{Beam Search 的适配}

标准 Beam Search 直接应用于 CTC 输出时需要特殊处理：

\begin{enumerate}
    \item CTC 输出中包含大量\textbf{blank token}，需要在 Beam Search 中正确处理前缀合并
    \item 需要集成\textbf{发音词典}（pronunciation dictionary），将音素序列映射到词
    \item 需要在每个时间步同时维护音素级和词级的假设
    \item 需要平衡解码速度和搜索质量——在 20ms 内完成
\end{enumerate}

\subsection{本章小结}

\begin{itemize}
    \item 语音 BCI 本质上是一个带有严格约束的 seq2seq 问题
    \item CTC 损失函数、GRU 编码器、Beam Search、语言模型——这些都是 NLP 领域的经典技术
    \item 但 BCI 的数据量和延迟约束要求对这些技术进行\textbf{针对性适配}
    \item 语言模型提供的先验知识对于弥补神经解码器的错误至关重要
\end{itemize}

%% ========== 总结与延伸 ==========

